% Part: incompleteness
% Chapter: representability-in-q
% Section: minimization-representable

\documentclass[../../../include/open-logic-section]{subfiles}

\begin{document}

\olfileid{inc}{req}{min}
\olsection{નિયમિત લઘુત્તમીકરણ~$\Th{Q}$માં નિરૂપ્ય છે}

હવે અમર્યાદિત શોધનો વિચાર કરીએ.
માનો કે $g(x, z)$ નિયમિત છે અને
!!{formula}~$!A_g(x, z, y)$ વડે
$\Th{Q}$માં નિરૂપ્ય છે.
માનો કે $f$ને
$f(z) = \umin{x}{[g(x, z) = 0]}$
દ્વારા વ્યાખ્યાયિત કરીએ.
$f$ને નિરૂપતું
!!a{formula}~$!A_f(z, y)$ શોધવું છે.
$f(z)$નું મૂલ્ય એવી સંખ્યા~$x$
છે જે (a) $g(x, z) = 0$ સંતોષે
છે અને (b) એવી સંખ્યાઓમાં સૌથી
નાની છે, એટલે દરેક $w < x$
માટે $g(w, z) \neq 0$ છે.
તેથી આ સ્વાભાવિક પસંદગી છે:
\[
!A_f(z,y) \ident !A_g(y, z, \Obj 0) \land \lforall[w][(w < y \lif \lnot
  !A_g(w, z, \Obj 0))].
\]
સામાન્ય કિસ્સામાં, અલબત્ત,
$z$ને $z_0$, \dots, $z_k$થી
બદલવું પડશે.

સાબિતીમાં ફરીથી એ બાબતનાં
કેટલાંક લેમ્મા આવશે કે
$\Th{Q}$ શું સાબિત કરવા જેટલું
સમર્થ છે.

\begin{lem}
\ollabel{lem:succ} દરેક
!!{constant}~$a$ અને દરેક
પ્રાકૃતિક સંખ્યા~$n$ માટે,
\[
\Th{Q} \Proves \eq[(a' + \num n)][(a + \num n)'].
\]
\end{lem}

\begin{proof}
હંમેશની જેમ, $n$ પર અનુમાન
વડે સાબિતી કરીએ. આરંભિક કિસ્સામાં
$n = 0$ છે; એટલે $\Th{Q}$માં
$\eq[(a' + \Obj 0)][(a + \Obj
0)']$ સાબિત કરવાનું છે.
આપણી પાસે આ છે:
\begin{align}
  \Th{Q} & \Proves \eq[(a' + \Obj 0)][a'] \quad \text{સ્વયંસિદ્ધિ $Q_4$ વડે}
  \ollabel{step1}\\
  \Th{Q} & \Proves  \eq[(a + \Obj 0)][a] \quad \text{સ્વયંસિદ્ધિ $Q_4$ વડે}
  \ollabel{step2} \\
  \Th{Q} & \Proves \eq[(a + \Obj 0)'][a'] \quad \text{\olref{step2} વડે}
  \ollabel{step3} \\
  \Th{Q} & \Proves \eq[(a' + \Obj 0)][(a + \Obj 0)'] \quad
  \text{\olref{step1} અને \olref{step3} વડે}\notag
\end{align}
અનુમાનના પગલામાં ધારો કે
$\Th{Q} \Proves \eq[(a' + \num n)][(a + \num n)']$
બતાવી ચૂક્યા છીએ. કારણ કે
$\num{n+1}$ એટલે $\num{n}'$,
હવે $\Th{Q}$માં
$\eq[(a' + \num{n}')][(a + \num{n}')']$
સાબિત કરવાનું છે. આપણે મેળવીએ:
\begin{align}
  \Th{Q} & \Proves \eq[(a' + \num n')][(a' + \num n)'] \quad
  \text{સ્વયંસિદ્ધિ $!Q_5$ વડે} \ollabel{step5}\\
  \Th{Q} & \Proves \eq[(a' + \num n)'][(a + \num n')'] \quad
  \text{અનુમાનની ધારણા અને સ્વયંસિદ્ધિ $!Q_5$ વડે} \ollabel{step6}\\
  \Th{Q} & \Proves \eq[(a' + \num n')][(a + \num n')'] \quad
  \text{\olref{step5} અને \olref{step6} વડે.} \notag
\end{align}
\sourcecorrection{OLINC-029}{સ્થિર અંગ્રેજી
મૂળના અનુમાન-પગલાની બીજી પંક્તિમાં
સાબિત કરવાનું અંતિમ સમીકરણ જ
અનુમાનની ધારણા તરીકે લખાયું છે;
છેલ્લી પંક્તિમાં મધ્યવર્તી સમીકરણ
આવી જાય છે. અહીં અનુમાનની
ધારણા તથા $Q_5$ વડે મળતું
મધ્યવર્તી સમીકરણ બીજી પંક્તિમાં
અને માંગેલું તારણ છેલ્લી પંક્તિમાં
આપ્યું છે.}
\end{proof}

ફરી નોંધવું જોઈએ કે આ દાવો
$\Th{Q}$ આ
$\lforall[x][\lforall[y][(x' + y) = (x + y)']]$
સાબિત કરે છે એવા દાવા કરતાં
નબળો છે. આ !!{sentence}
$\Struct{N}$માં સાચું હોવા છતાં
$\Th{Q}$ તેને સાબિત કરતું નથી.

\begin{lem}
\ollabel{lem:less-zero}
$\Th{Q} \Proves \lforall[x][\lnot x < \Obj 0]$.
\end{lem}

\begin{proof}
સાબિતી અનૌપચારિક રીતે આપીએ
(એટલે ઔપચારિક !!{derivation}
કેવી રીતે રચવું તેના માત્ર સંકેતો
આપીએ).

કોઈપણ~$a$ માટે $\lnot a < \Obj 0$
સાબિત કરવાનું છે. $<$ની
વ્યાખ્યા પ્રમાણે $\Th{Q}$માં
$\lnot \lexists[y][\eq[(y'+a)][\Obj 0]]$
સાબિત કરવાનું છે. આપણે
$\lexists[y][\eq[(y'+a)][\Obj 0]]$
ધારીને વિરોધાભાસ મેળવીશું.
માનો કે $\eq[(b'+a)][\Obj 0]$.
$!Q_3$ વડે
$\eq[a][\Obj 0] \lor \lexists[y][\eq[a][y']]$
મળે છે. હવે બે કિસ્સા જોઈએ.

કિસ્સો 1: $\eq[a][\Obj 0]$.
$\eq[(b'+a)][\Obj 0]$માંથી
$\eq[(b' + \Obj 0)][\Obj 0]$
મળે છે. $\Th{Q}$ની સ્વયંસિદ્ધિ
$!Q_4$ વડે
$\eq[(b' + \Obj 0)][b']$,
એટલે $\eq[b'][\Obj 0]$ મળે છે.
પણ સ્વયંસિદ્ધિ $!Q_2$ વડે
$\eq/[b'][\Obj 0]$ પણ મળે છે;
આ વિરોધાભાસ છે.

કિસ્સો 2: કોઈ~$c$ માટે
$\eq[a][c']$. તો
$\eq[(b' + c')][\Obj 0]$ મળે છે.
સ્વયંસિદ્ધિ $!Q_5$ વડે
$\eq[(b' + c)'][\Obj 0]$ મળે,
જે ફરી સ્વયંસિદ્ધિ $Q_2$ની
વિરુદ્ધ છે.
\end{proof}

\begin{lem}
\ollabel{lem:less-nsucc}
દરેક પ્રાકૃતિક સંખ્યા~$n$ માટે,
  \[
  \Th{Q} \Proves
  \lforall[x][(x < \num {n+1} \lif (\eq[x][\Obj 0] \lor \dots \lor
    \eq[x][\num n]))].
  \]
\end{lem}

\begin{proof}
$n$ પર અનુમાન વાપરીએ.
પહેલાં $n = 0$નો આરંભિક
કિસ્સો જોઈએ. ત્યારે કોઈપણ~$a$
માટે $a < \num 1 \lif
\eq[a][\Obj 0]$ બતાવવું છે.
માનો કે $a < \num 1$. તો
$<$ની વ્યાખ્યા આપતી સ્વયંસિદ્ધિ
પ્રમાણે
$\lexists[y][\eq[(y'+a)][\Obj 0']]$
છે (કારણ કે $\num 1 \ident
\Obj 0'$).

માનો કે $b$ આ ગુણધર્મ
ધરાવે છે, એટલે
$\eq[(b'+a)][\Obj 0']$.
$\eq[a][\Obj 0]$ બતાવવાનું છે.
સ્વયંસિદ્ધિ $!Q_3$ પ્રમાણે
$\eq[a][\Obj 0]$ અથવા
કોઈ~$c$ માટે $\eq[a][c']$.
પહેલા કિસ્સામાં બતાવવાનું કંઈ
બાકી નથી. તેથી
$\eq[a][c']$ ધારો. તો
$\eq[(b' + c')][\Obj 0']$.
$\Th{Q}$ની સ્વયંસિદ્ધિ $!Q_5$
વડે $\eq[(b'+c)'][\Obj 0']$.
સ્વયંસિદ્ધિ $!Q_1$ વડે
$\eq[(b' + c)][\Obj 0]$.
પણ સ્વયંસિદ્ધિ $!Q_8$ પ્રમાણે
આનો અર્થ $c < \Obj 0$ છે,
જે \olref{lem:less-zero}થી
મળેલા પરિણામની વિરુદ્ધ છે.

હવે અનુમાનનું પગલું જોઈએ.
$n$ માટેનો કિસ્સો ધારીને
$n+1$ માટે સાબિત કરીએ.
માનો કે $a < \num {n+2}$.
ફરી $!Q_3$ વડે બે કિસ્સા:
$\eq[a][\Obj 0]$ અથવા
કોઈ~$b$ માટે $\eq[a][b']$.
પહેલા કિસ્સામાં
$\eq[a][\Obj 0] \lor \dots \lor
\eq[a][\num{n+1}]$ તરત મળે છે.
બીજા કિસ્સામાં
$b' < \num {n+2}$, એટલે
$b' < \num{n+1}'$.
સ્વયંસિદ્ધિ $!Q_8$ વડે કોઈ~$c$
માટે
$\eq[(c'+b')][\num{n+1}']$.
સ્વયંસિદ્ધિ $!Q_5$ વડે
$\eq[(c'+b)'][\num{n+1}']$.
સ્વયંસિદ્ધિ $!Q_1$ વડે
$\eq[(c'+b)][\num{n+1}]$,
અને તેથી સ્વયંસિદ્ધિ $!Q_8$
વડે $b < \num{n+1}$.
અનુમાનની ધારણા પ્રમાણે
$\eq[b][\Obj 0] \lor \dots \lor
\eq[b][\num{n}]$. તર્કથી
$\eq[b'][\Obj 0'] \lor \dots \lor
\eq[b'][\num{n}']$ મળે છે;
અને $\eq[a][b']$ હોવાથી
$\eq[a][\num{1}] \lor \dots \lor
\eq[a][\num{n+1}]$ મળે છે.
\end{proof}

\begin{lem}
  \ollabel{lem:trichotomy} દરેક પ્રાકૃતિક સંખ્યા~$m$ માટે,
  \[
  \Th{Q} \Proves
  \lforall[y][((y < \num{m} \lor \num{m} < y) \lor \eq[y][\num{m}])].
  \]
\end{lem}

\begin{proof}
$m$ પર અનુમાન કરીએ.
પહેલાં $m=0$નો કિસ્સો લો.
$!Q_3$ વડે
$\Th{Q} \Proves
\lforall[y][(\eq[y][\Obj 0] \lor \lexists[z][\eq[y][z']])]$.
કોઈપણ~$a$ લો. તો
$\eq[a][\Obj 0]$ અથવા
કોઈ~$b$ માટે $\eq[a][b']$.
પહેલા કિસ્સામાં
$(a < \Obj 0 \lor \Obj
0 < a) \lor \eq[a][\Obj 0]$
પણ મળે છે. જો $\eq[a][b']$,
તો~$\eq$ માટેના તર્ક વડે
$\eq[(b' + \Obj 0)][(a + \Obj 0)]$.
$!Q_4$ વડે
$\eq[(a + \Obj 0)][a]$;
તેથી $\eq[(b' + \Obj 0)][a]$,
અને પરિણામે
$\lexists[z][\eq[(z' + \Obj 0)][a]]$.
$!Q_8$ની $<$ની વ્યાખ્યા
પ્રમાણે $\Obj 0 < a$.
જો $\Obj 0 < a$, તો
$(\Obj 0 < a \lor a
< \Obj 0) \lor \eq[a][\Obj 0]$
પણ મળે છે.

હવે ધારો કે આ સાબિત થયું છે:
\begin{align*}
  \Th{Q} & \Proves \lforall[y][((y < \num{m} \lor \num{m} < y) \lor
    \eq[y][\num{m}])]
  \intertext{અને આપણે આ બતાવવું છે:}
  \Th{Q} & \Proves \lforall[y][((y < \num{m+1} \lor \num{m+1} < y) \lor
    \eq[y][\num{m+1}])]
\end{align*}
કોઈપણ~$a$ લો. $!Q_3$ પ્રમાણે
$\eq[a][\Obj 0]$ અથવા
કોઈ~$b$ માટે $\eq[a][b']$.
પહેલા કિસ્સામાં $!Q_4$ વડે
$\eq[\num{m}' +
a][\num{m+1}]$; એટલે $!Q_8$
વડે $a < \num{m+1}$.

હવે બીજો કિસ્સો લો:
$\eq[a][b']$. અનુમાનની
ધારણા પ્રમાણે
$(b < \num{m} \lor \num{m} < b) \lor
    \eq[b][\num{m}]$.

પહેલો વિકલ્પ $b < \num{m}$
$!Q_8$ પ્રમાણે
$\lexists[z][\eq[(z' + b)][\num{m}]]$
ની સમકક્ષ છે. માનો કે~$c$ આ
ગુણધર્મ ધરાવે છે. જો
$\eq[(c' + b)][\num{m}]$,
તો $\eq[(c' + b)'][\num{m}']$
પણ મળે છે. $!Q_5$ વડે
$\eq[(c' + b')][(c' + b)']$.
એટલે $\eq[(c' +
b')][\num{m}']$. હવે~$c$
પર અસ્તિત્વ-પરિમાણક લાવીને
અને $\num{m}' \ident
\num{m+1}$ ધ્યાનમાં રાખીને
$\lexists[u][\eq[(u' + b')][\num{m+1}]]$
મળે છે. તેથી $b < \num{m}$
હોય, તો $b' < \num{m+1}$;
અને $a < \num{m+1}$.
\sourcecorrection{OLINC-030}{સ્થિર અંગ્રેજી
મૂળ અહીં $Q_5$ વડે નિરર્થક
સ્વસમાનતા લખે છે અને પછી
$c'$ પર અસ્તિત્વ-પરિમાણક
લાવવાનું કહે છે. જરૂરી સમીકરણ
$(c'+b')=(c'+b)'$ છે;
તેમાં સાક્ષી~$c$ જ છે, જેને
પરિમાણિત કરતાં $b'<\num{m+1}$
મળે છે.}

હવે $\num{m} < b$ માનો,
એટલે
$\lexists[z][\eq[(z' +
\num{m})][b]]$. માનો કે~$c$
આવું~$z$ છે, એટલે
$\eq[(c' + \num{m})][b]$.
તર્કથી
$\eq[(c' + \num{m})'][b']$.
$!Q_5$ વડે
$\eq[(c' + \num{m}')][b']$.
$\eq[a][b']$ અને
$\num{m}' \ident \num{m+1}$
હોવાથી
$\eq[(c' + \num{m+1})][a]$.
$!Q_8$ વડે $\num{m+1} < a$.

છેવટે $\eq[b][\num{m}]$
ધારો. તો તર્કથી
$\eq[b'][\num{m}']$;
એટલે $\eq[a][\num{m+1}]$.

આથી $m$ અને~$b$ માટેના
દરેક વિકલ્પમાંથી $m+1$
અને~$a$ માટેનો અનુરૂપ
વિકલ્પ મળે છે.
\end{proof}

\begin{prop}
\ollabel{prop:rep-minimization}
$!A_g(x, z, y)$ એ~$g(x, z)$ને
$\Th{Q}$માં નિરૂપે, તો
\[
!A_f(z,y) \ident !A_g(y, z, \Obj 0) \land \lforall[w][(w < y \lif \lnot
  !A_g(w, z, \Obj 0))]
\]
એ $f(z) = \umin{x}{[g(x, z) = 0]}$ને
નિરૂપે છે.
\end{prop}

\begin{proof}
પહેલાં બતાવીએ કે $f(n) = m$
હોય, તો
$\Th{Q} \Proves !A_f(\num n, \num m)$,
એટલે કે
\begin{align}
  \Th{Q} & \Proves !A_g(\num{m}, \num{n}, \Obj 0) \land \lforall[w][(w <
    \num{m} \lif \lnot !A_g(w, \num{n}, \Obj 0))]. \notag
  \intertext{કારણ કે $!A_g(x, z, y)$ એ $g(x, z)$ને નિરૂપે છે અને $f(n) = m$ હોય તો $g(m, n) = 0$, તેથી}
\Th{Q} & \Proves !A_g(\num{m}, \num{n}, \Obj 0). \notag
\intertext{$f(n) = m$ હોય, તો દરેક $k < m$ માટે $g(k, n) \neq 0$. તેથી}
\Th{Q} & \Proves \lnot !A_g(\num{k}, \num{n}, \Obj 0). \notag
\intertext{આ પરથી મળે છે:}
\Th{Q} & \Proves \lforall[w][(w < \num{m} \lif \lnot
  !A_g(w, \num{n}, \Obj 0))]. \ollabel{rep-less}
\end{align}
$m = 0$ હોય તો
\olref{lem:less-zero} વડે અને
અન્યથા \olref{lem:less-nsucc}
વડે છેલ્લું તારણ મળે છે.

હવે $f(n) = m$ હોય, તો
$\Th{Q} \Proves
\lforall[y][(!A_f(\num{n}, y) \lif
\eq[y][\num{m}])]$
તે બતાવીએ. દલીલ ફરી
અનૌપચારિક રીતે દર્શાવીએ અને
ઔપચારિક રૂપાંતર વાચક પર છોડીએ.

માનો કે $!A_f(\num{n}, b)$.
ત્યારે (a) $!A_g(b, \num{n},
\Obj 0)$ અને (b)
$\lforall[w][(w < b \lif \lnot !A_g(w, \num{n}, \Obj
  0))]$ મળે છે.
\olref{lem:trichotomy} પ્રમાણે
$(b < \num{m} \lor \num{m} < b)
\lor \eq[b][\num{m}]$.
$b < \num{m}$ અને $\num{m}
< b$ બંને વિરોધાભાસ આપે છે
તે બતાવીએ.

જો $\num{m} < b$, તો~(b)
પરથી $\lnot !A_g(\num{m},
\num{n}, \Obj 0)$ મળે છે.
પરંતુ $m = f(n)$ હોવાથી
$g(m, n) = 0$; અને~$!A_g$
એ~$g$ને નિરૂપે છે, એટલે
$\Th{Q} \Proves !A_g(\num{m},
\num{n}, \Obj 0)$.
આ વિરોધાભાસ છે.

હવે $b < \num{m}$ માનો.
\olref{rep-less} પ્રમાણે
$\Th{Q} \Proves \lforall[w][(w <
  \num{m} \lif \lnot !A_g(w, \num{n}, \Obj 0))]$;
તેથી $\lnot !A_g(b, \num{n},
\Obj 0)$ મળે છે. આ ફરી~(a)
ની વિરુદ્ધ છે.
\end{proof}

\end{document}
